\begin{Exercise}[label={ex:temperatura}]
    Una stazione meteorologica registra la temperatura esterna ogni ora, ogni giorno per un mese intero, in una struttura dati apposita: una matrice bi-dimensionale di dimensioni fisiche di 31 righe (i giorni del mese) per 24 colonne (le ore del giorno).
    Si scriva un programma che calcoli i seguenti dati:
    \begin{itemize}
        \item Temperatura media del mese.
        \item Temperatura media giornaliera.
        \item Temperatura media diurna (ore comprese tra le 7:00 e le 19:00), su tutto il mese.
        \item Temperatura media notturna su tutto il mese.
    \end{itemize}

    \textbf{Nota: } non tutti i mesi hanno 31 giorni! Popoliamo la matrice utilizzando la funzione \texttt{rand()} della libreria \texttt{stdlib}.
\end{Exercise}

\begin{Answer}[ref={ex:temperatura}]
  	\lstinputlisting[style=cstylesmall]{src/temperatura.c}
\end{Answer}
